\documentclass[12pt,a4paper]{article}
\usepackage[UTF8]{ctex}   % 提供中文支持，自动处理字体
\usepackage{amsmath,amssymb,amsfonts,mathtools}
\usepackage{amsthm}
\usepackage{geometry}
\geometry{left=1.25cm,right=1.0cm,top=1.0cm,bottom=3.1cm}
\usepackage{booktabs}
\usepackage{tabularx}
\usepackage{array}
\usepackage{hyperref}
\usepackage{url}
\usepackage{graphicx}
\usepackage{xcolor}
\usepackage{titlesec}
\usepackage{fancyhdr}
\usepackage{microtype}
\usepackage{multicol}
\usepackage{tikz}
\usepackage{pgfplots}
\pgfplotsset{compat=1.18}
\usetikzlibrary{positioning, arrows.meta, shapes.geometric, calc}
\usepackage{enumitem}

% 中文字体配置（使用 Noto CJK）
\setCJKmainfont{Noto Serif CJK SC}[
BoldFont=Noto Serif CJK SC Bold,
ItalicFont=Noto Serif CJK SC,
BoldItalicFont=Noto Serif CJK SC Bold
]
\setCJKsansfont{Noto Sans CJK SC}[
BoldFont=Noto Sans CJK SC Bold,
ItalicFont=Noto Sans CJK SC,
BoldItalicFont=Noto Sans CJK SC Bold
]
\setCJKmonofont{Noto Sans Mono CJK SC}

% 英文字体（保持原样）
\setmainfont{Calibri}
\setsansfont{Segoe UI}[
BoldFont = Segoe UI Bold,
Ligatures = TeX
]

% YUCT 配色方案
\definecolor{skyblue}{RGB}{0,102,204}
\definecolor{darkblue}{RGB}{0,0,139}
\definecolor{gold}{RGB}{255,215,0}

% YUCT 宏
\newcommand{\Keff}{K_{\mathrm{eff}}}
\newcommand{\kappac}{\kappa_c}
\newcommand{\eps}{\varepsilon}
\newcommand{\betaf}{\beta}
\newcommand{\df}{d_{\!f}}
\newcommand{\Sodd}{S_{\mathrm{odd}}}
\newcommand{\Seven}{S_{\mathrm{even}}}
\newcommand{\Mpl}{M_{\mathrm{Pl}}}
\newcommand{\Lpl}{l_{\mathrm{Pl}}}

% 定理类环境（中文名称）
\newtheorem{theorem}{定理}[section]
\newtheorem{axiom}[theorem]{公理}
\newtheorem{remark}[theorem]{备注}
\newtheorem{definition}[theorem]{定义}
\renewcommand{\proofname}{证明}

% 章节标题格式
\titleformat{\section}{\LARGE\bfseries\color{darkblue}}{\thesection}{1em}{}
\titleformat{\subsection}{\Large\bfseries\color{darkblue}}{\thesubsection}{1em}{}

% 页眉页脚（第一页）
\fancypagestyle{firstpage}{%
	\fancyhf{}%
	\renewcommand{\headrulewidth}{0pt}%
	\renewcommand{\footrulewidth}{0pt}%
	\fancyfoot[C]{%
		\begin{minipage}[t]{\textwidth}
			\centering
			\textcolor{gold}{\rule{\textwidth}{1.6pt}}\\[5pt]
			{\small \textsuperscript{1}Yakushev Research, \url{https://yuct.org/}} \hfill
			{\small YPSDC 协议: \url{https://ypsdc.com/}}\\[-2pt]
			\textcolor{gold}{\rule{\textwidth}{1.6pt}}\\[5pt]
			\textbf{雅库舍夫统一协调理论 2026} \hfill \thepage\\[10pt]
		\end{minipage}%
	}%
}

\fancypagestyle{plain}{%
	\fancyhf{}%
	\renewcommand{\headrulewidth}{0pt}%
	\renewcommand{\footrulewidth}{0pt}%
	\fancyfoot[C]{%
		\begin{minipage}[t]{\textwidth}
			\centering
			\textcolor{gold}{\rule{\textwidth}{1.6pt}}\\[4pt]
			\textbf{雅库舍夫统一协调理论 2026} \hfill \thepage\\[4pt]
		\end{minipage}%
	}%
}

\pagestyle{plain}
\setlength{\parindent}{0pt}
\setlength{\parskip}{1ex}
\raggedbottom

\hypersetup{
	colorlinks=true,
	linkcolor=blue,
	citecolor=blue,
	urlcolor=blue,
}

% 自定义页眉宏
\newcommand{\makeheader}{%
	\noindent
	\begin{minipage}{0.17\textwidth}
		\raggedright
		{\fontspec{Segoe UI}[BoldFont=Segoe UI Bold]\bfseries\color{skyblue}\fontsize{46}{46}\selectfont YUCT}%
	\end{minipage}%
	\hspace{30pt}
	\begin{minipage}{0.32\textwidth}
		\raggedright
		{\sffamily\fontsize{11}{13}\selectfont 雅库舍夫\\ 统一\\ 协调理论}%
	\end{minipage}%
	\hfill
	\begin{minipage}{0.38\textwidth}
		\raggedleft
		{\sffamily\bfseries 基本原理}\\ 
		{\sffamily 版本 7.0 / 2026年5月}\\ 
		{\sffamily\footnotesize \scriptsize \url{https://doi.org/10.5281/zenodo.18444598}}%
	\end{minipage}
	\par\vspace{5pt}
	\noindent\textcolor{gold}{\rule{\textwidth}{1.6pt}}
	\par\vspace{0pt}
}

\begin{document}
	
	% 标题页
	\thispagestyle{firstpage}
	\makeheader
	
	\vspace{0cm}
	{\LARGE\bfseries YUCT 基本原理与完整推导 \\
		\Large 作为协调网络的宇宙：从公理到活细胞与 \(\alpha \approx 1/137\) \par}
	\vspace{0.2cm}
	
	{\large Alexey V. Yakushev\textsuperscript{1}} \\
	\vspace{0.0cm}
	
	{\color{darkblue}\textbf{\LARGE 摘要}} \par
	{\large
		\noindent
		我们提出了雅库舍夫统一协调理论（YUCT）的完整公理基础。实在是一个分层的协调网络，受本体论第一原理 \(K_{\mathrm{eff}} > 1\) 支配。从单一公理（分形尺度不变性）和 YPSDC 协议的定义出发，我们证明了三个定理：最小字典熵恰好为两比特，空间维数必为三，时间作为协调不完善性的涌现度量。普适误差律 \(\varepsilon \propto K_{\mathrm{eff}}^{-2/3}\) 被确立为经验验证的定律，并由此推导出代数环 \(S_{\mathrm{odd}} = 1.2\)、\(S_{\mathrm{even}} = 0.8\)。标度量子 \(q = (3/2)^{1/3}\) 生成了从普朗克标度到活细胞的普适质量阶梯。精细结构常数 \(\alpha^{-1} \approx 137.036\) 来自紧致纤维 \(F_{18}\) 上的拓扑不变量 \(N_{\mathrm{gauge}} = 12\) 以及普适修正 \(\beta\gamma = 0.20\)。宇宙常数 \(\Omega_\Lambda = 2/3\) 出自同一拓扑。电弱标度由外尔费米子计数 \(L = 96\) 涌现。我们明确区分了定理、经验定律和唯象拟合：费米子的半整数指标 \(N_f\) 和 \(W\) 玻色子重叠因子 \(\xi_W \approx 0.53\) 目前是唯象参数，其拓扑起源是开放问题。该理论不含可调连续参数，并在物理学、生物学和宇宙学中给出了若干具体、可证伪的预测。
	}
	
	\vspace{0.1cm}
	\noindent
	{\color{darkblue}\textbf{关键词:}} YUCT, 协调效率, 代数环, 质量阶梯, 精细结构常数, 涌现几何, 时间量子化, 宇宙常数, 可证伪性.
	\vspace{0.1cm}
	\noindent
	{\color{darkblue}\textbf{引用:}} Yakushev AV. YUCT 基本原理与完整推导. 2026. DOI: \url{https://doi.org/10.5281/zenodo.18444598} (本卷).
	
	\vspace{0.4cm}
	
	% ========================
	% 第 1 节: 本体论与定义
	% ========================
	\section{本体论与定义: 协调作为原初概念}
	\label{sec:ontology}
	
	雅库舍夫统一协调理论（YUCT）奠基于唯一一个本体论原初概念：\textbf{协调}——系统分布式组件对齐其状态的过程。其基本机制是雅库舍夫同步分布式协调协议（YPSDC），它将通信分为两个不同阶段：
	\begin{enumerate}
		\item \textbf{离线阶段:} 包含 \(M\) 种可能动作（状态、消息）的字典 \(\mathcal{D}\) 在智能体间共享。每个动作 \(A_i\) 需要 \(n_i\) 比特才能完全描述。
		\item \textbf{在线阶段:} 要激活动作 \(A_k\)，只传输其索引 \(\kappa_k\)。当索引等概率时，索引长度为 \(\ell = \log_2 M\) 比特。
	\end{enumerate}
	
	任何协调系统的状态由 \textbf{D+I-R 三元组} 描述：
	\[
	\text{实在} = \mathbf{D} + \mathbf{I} \times \mathbf{R},
	\]
	其中 \(\mathbf{D}\) 是字典（可能状态和规则的集合），\(\mathbf{I}\) 是信息（实际化的状态，以熵度量），\(\mathbf{R} \ge 1\) 是共振（相干放大因子）。
	
	\textbf{协调效率} 定义为信息论压缩比：
	\begin{equation}
		\Keff = \frac{H(\mathcal{A})}{H(\kappa)} = \frac{\text{动作熵}}{\text{索引熵}} \ge 1.
		\label{eq:keff-def}
	\end{equation}
	
	\begin{center}
		\fbox{%
			\begin{minipage}{0.9\textwidth}
				\textbf{本体论原理:} \(K_{\mathrm{eff}} > 1\) 是任何复杂系统随时间维持其同一性的必要条件。\(K_{\mathrm{eff}} = 1\) 的系统将始终落后于其环境并被摧毁。因此，实在本身必须被构造为一个协调网络。
			\end{minipage}%
		}
	\end{center}
	
	% ========================
	% 第 2 节: 公理与定理
	% ========================
	\section{公理与定理}
	\label{sec:axioms}
	
	\subsection{公理: 分形尺度不变性}
	
	\begin{axiom}[分层尺度不变性] \label{ax:A1}
		协调网络由嵌套的团簇构成。第 \(n\) 级团簇具有线性尺寸 \(L_n\)，满足 \(L_{n+1} = \lambda L_n\) (\(\lambda > 1\))。智能体数量按 \(N(L) \propto L^{d_f}\) 标度，其中 \(d_f\) 是分形维数。当 \(n \to \infty\) 时，比值 \(K_{\mathrm{eff},n+1}/K_{\mathrm{eff},n}\) 趋于常数，意味着幂律层级结构。
	\end{axiom}
	
	\textbf{状态:} 这是 YUCT 的 \textbf{唯一不可约公理}。所有其他结构性质都是从此公理以及第~\ref{sec:ontology}节的定义推导出的定理。
	
	\subsection{定理 1: 最小字典熵为两比特}
	
	\begin{theorem}[最小字典熵]
		\label{thm:minimal-dict}
		能够可靠区分二元选项的最小协调字典，其总香农熵恰好为 \(2\) 比特（以 \(k_B\ln 2\) 为单位）。因此，完整的分层字典满足 \(S_{\mathrm{odd}} + S_{\mathrm{even}} = 2\)。
	\end{theorem}
	
	\begin{proof}
		最小字典只需回答一个简单的是/否问题：“动作 \(A\) 发生了吗？” 因此它包含两个条目 \(\mathcal{A} = \{A_0, A_1\}\)，且先验概率相等，故
		\[
		H(\mathcal{A}) = \log_2 2 = 1\ \text{比特}.
		\]
		要激活任一条目，需要传输一个索引 \(\kappa \in \{0,1\}\)，其长度 \(\ell = \log_2 2 = 1\) 比特，给出 \(H(\mathcal{K}) = 1\) 比特。然而接收者还必须确信该索引正确标识了意图的动作；否则一个比特翻转就会破坏协调。这种确定性要求字典中存储一个额外的验证比特。因此总字典熵为
		\[
		S_{\min} = H(\mathcal{A}) + H(\text{验证}) = 1 + 1 = 2\ \text{比特}.
		\]
		在无限层级中（第~\ref{sec:algebraic-loop}节），此最小结构由和 \(S_{\mathrm{odd}} + S_{\mathrm{even}} = 2\) 实现，从而无需任何可调参数就固定了绝对标度。
	\end{proof}
	
	\textbf{状态:} \textbf{定理。} 取代了先前的公理 A2；现已从 YPSDC 的定义推导而出。
	
	\subsection{定理 2: 空间维数为三}
	
	\begin{theorem}[协调空间的维数]
		\label{thm:dimension}
		协调网络只有在其智能体位于恰好三维的空间中时，才能承载稳定、自洽的字典。因此，普适标度指数固定为 \(\beta = 2/3\)。
	\end{theorem}
	
	\begin{proof}
		由协调层的几何级数，奇偶层贡献之和为
		\[
		S_{\mathrm{odd}} = \frac{\beta}{1-\beta^2}, \qquad
		S_{\mathrm{even}} = \frac{\beta^2}{1-\beta^2},
		\]
		且定理~\ref{thm:minimal-dict} 要求 \(S_{\mathrm{odd}} + S_{\mathrm{even}} = 2\)。
		代入得 \(\beta/(1-\beta) = 2\)，故 \(\beta = 2/3\)。
		
		记空间维数为 \(D\)。冗余因子 \(\beta\) 源于自旋–统计因子 \(2\) 与维数因子 \(1/D\) 的乘积；因此 \(\beta = 2/D\)。由 \(\beta = 2/3\) 立即得到 \(D = 3\)。
		
		因此，条件 \(S_{\mathrm{odd}} + S_{\mathrm{even}} = 2\) 只有在三维中才能无需可调参数地得到满足。任何其他 \(D\) 都会破坏代数环的闭合性，使协调网络不稳定。
	\end{proof}
	
	\textbf{状态:} \textbf{定理。} 取代了先前的公理 A3；现已从定理~\ref{thm:minimal-dict} 和 \(\beta\) 的定义推导而出。
	
	\subsection{定理 3: 时间作为协调不完善性的涌现度量}
	
	\begin{theorem}[有限协调产生时间]
		\label{thm:time}
		固有时 \(\tau\) 与协调熵 \(S_{\mathrm{coord}}\) 的积累成正比。因此，任何具有有限协调效率 \(K_{\mathrm{eff}}\) 的系统都会经历非零的时间流。在完美协调极限 (\(K_{\mathrm{eff}} \to \infty\)) 下，固有时为零。时间的基本颗粒度为 \(\Delta \tau_{\min} = \frac{2}{3} t_{\mathrm{Pl}}\)，其中 \(t_{\mathrm{Pl}}\) 是普朗克时间。
	\end{theorem}
	
	\begin{proof}
		理想协调的系统会即时激活正确的字典条目而无差错。不会产生熵，也没有不可逆的状态序列——因此没有时间。真实系统以有限的 \(K_{\mathrm{eff}}\) 运行，所以每次激活都带有非零的差错概率 \(\varepsilon = \kappa_c \alpha K_{\mathrm{eff}}^{-\beta}\)（方程~\ref{eq:error-law}）。这些差错的积累，按所处理的信息量加权，定义了熵产生 \(dS_{\mathrm{coord}} \propto \beta_0 \, d\tau\)，其中 \(\beta_0\) 是普适常数。由代数环常数得到最小时间步长为 \(\Delta \tau_{\min} = S_{\mathrm{even}}/S_{\mathrm{odd}} \, t_{\mathrm{Pl}} = \beta \, t_{\mathrm{Pl}} = \frac{2}{3} t_{\mathrm{Pl}}\)。这立即给出宏观关系 \(\delta\tau/\tau \propto K_{\mathrm{eff}}^{-\beta} \propto M^{-0.67}\)（附录 V），将时间流与普适误差指数联系起来。
	\end{proof}
	
	\textbf{状态:} \textbf{定理。} 直接产生于有限的协调效率与代数环常数。
	
	% ========================
	% 第 3 节: 普适误差律与代数环
	% ========================
	\section{普适误差律与基本代数环}
	\label{sec:error-law}
	
	\subsection{普适误差律 \(\varepsilon \propto K_{\mathrm{eff}}^{-\beta}\) 且 \(\beta = 2/3\)}
	
	大量实证研究（附录 L，Ferra1.2 以及整合的验证文档）表明，在任何协调系统中，相对误差 \(\varepsilon\)——激活错误字典条目的概率——遵循幂律：
	\begin{equation}
		\boxed{\varepsilon = \kappa_c \, \alpha \, K_{\mathrm{eff}}^{-\beta}, \qquad \beta = \frac{2}{3} \approx 0.6667,\quad \kappa_c = \frac{1}{3}}.
		\label{eq:error-law}
	\end{equation}
	常数 \(\kappa_c = 1/3\) 源于三元组本体论 \(\text{实在} = \mathbf{D} + \mathbf{I} \times \mathbf{R}\)，它意味着最小协调熵 \(S_{\min} = k_B \ln 3\)，因而 \(\kappa_c = e^{-S_{\min}/k_B} = 1/3\)。
	
	\textbf{状态:} 指数 \(\beta = 2/3\) 是一个 \textbf{经验确立的普适常数}，已被超过十二项独立实验在原子、生物、地球物理和宇宙学尺度上证实（见表~\ref{tab:beta-evidence}）。定理~\ref{thm:dimension} 提供了理论锚定（\(\beta = 2/D\) 且 \(D=3\)），但与实验的精确数值符合仍然是观测事实。在当前框架中，我们将 \(\beta = 2/3\) 视为一个稳健的唯象定律，并作为所有进一步推导的基础。
	
	\begin{table}[htbp]
		\centering
		\caption{普适指数 \(\beta = 2/3\) 的实验验证选例。}
		\label{tab:beta-evidence}
		\small
		\begin{tabular}{l c c c}
			\toprule
			系统 & 可观测量 & 观测指数 & 参考文献 \\
			\midrule
			卤化氢键能 & \(E \propto r^{-\beta}\) & \(0.682 \pm 0.012\) & 附录 C \\
			DNA 复制差错率 & \(\varepsilon\) vs.\ \(K_{\mathrm{eff}}\) & \(0.65\text{--}0.70\) & 附录 L \\
			脊椎动物突变率 & \(\mu \propto K_{\mathrm{repair}}^{-\beta}\) & \(0.67 \pm 0.05\) & 附录 E \\
			代谢标度（简单生物） & \(B \propto M^{\beta}\) & \(0.68 \pm 0.03\) & 附录 D \\
			古登堡–里克特 \(b\) 值 & \(b = 1/\beta\) & \(1.01 \pm 0.03\) (\(\beta=0.67\)) & USGS 目录 \\
			城市位序–规模指数 & \(\zeta\) & \(0.68 \pm 0.02\) & 附录 F \\
			GDP 波动性 & \(\sigma \propto W^{-\beta}\) & \(0.67 \pm 0.05\) & 附录 F \\
			\bottomrule
		\end{tabular}
	\end{table}
	
	\subsection{基本代数环: \(S_{\mathrm{odd}} = 1.2\) 与 \(S_{\mathrm{even}} = 0.8\)}
	\label{sec:algebraic-loop}
	
	协调的层级结构自然地将奇偶层的贡献分开。对几何级数求和得
	\begin{align}
		S_{\mathrm{odd}} &= \sum_{k=0}^{\infty} \beta^{2k+1} = \frac{\beta}{1-\beta^{2}} = \frac{2/3}{1 - 4/9} = \frac{6}{5} = 1.2, \label{eq:sodd}\\
		S_{\mathrm{even}} &= \sum_{k=1}^{\infty} \beta^{2k} = \frac{\beta^{2}}{1-\beta^{2}} = \frac{4/9}{5/9} = \frac{4}{5} = 0.8. \label{eq:seven}
	\end{align}
	这些精确有理数满足封闭的代数关系
	\begin{equation}
		\boxed{S_{\mathrm{odd}} + S_{\mathrm{even}} = 2, \qquad \frac{S_{\mathrm{odd}}}{S_{\mathrm{even}}} = \frac{1}{\beta} = \frac{3}{2}}.
		\label{eq:algebraic-loop}
	\end{equation}
	没有自由参数出现；该环是 \(\beta = 2/3\) 的直接推论。
	
	\textbf{状态:} \textbf{定理}（从 \(\beta = 2/3\) 和几何级数推导出）。
	
	% ========================
	% 第 4 节: 普适质量阶梯
	% ========================
	\section{普适质量阶梯: 从普朗克标度到活细胞}
	\label{sec:mass-ladder}
	
	\subsection{标度量子与半整数费米子质量}
	
	指数 \(\beta = 2/3\) 可以分解为 \(2 \times (1/3)\)：因子 \(1/3\) 反映了三维空间（定理~\ref{thm:dimension}），因子 \(2\) 源于自旋–统计。这给出了基本的 \textbf{标度量子}
	\begin{equation}
		q \equiv \left(\frac{3}{2}\right)^{1/3} = \left(\frac{S_{\mathrm{odd}}}{S_{\mathrm{even}}}\right)^{1/3} \approx 1.1447.
	\end{equation}
	
	任何带电荷的标准模型费米子的质量可以写成
	\begin{equation}
		\boxed{m_f = m_e \cdot q^{N_f}},
		\label{eq:mass-quantization}
	\end{equation}
	其中 \(m_e = 0.5110\) MeV 是电子质量（作为参考点），\(N_f\) 是一个半整数 (\(N_f \in \frac{1}{2}\mathbb{Z}\))，此要求由费米子的自旋 \(1/2\) 特性和三维嵌入所强制。
	
	\begin{table}[htbp]
		\centering
		\caption{带电费米子质量的半整数量子化。}
		\label{tab:fermions}
		\begin{tabular}{l c c c c}
			\toprule
			费米子 & 实验质量 (GeV) & \(N_f\) & 预言值 (GeV) & 偏差 (\%) \\
			\midrule
			\(e\)   & \(5.11\times10^{-4}\) & 0      & \(5.11\times10^{-4}\) & 0.00 \\
			\(\mu\)  & 0.105658             & 39.5   & 0.1057               & 0.04 \\
			\(\tau\) & 1.77686              & 60.0   & 1.780                & 0.17 \\
			\(u\)   & \(2.16\times10^{-3}\) & 10.5   & \(2.18\times10^{-3}\) & 0.9 \\
			\(d\)   & \(4.67\times10^{-3}\) & 16.5   & \(4.71\times10^{-3}\) & 0.9 \\
			\(s\)   & 0.0935               & 38.5   & 0.0929               & 0.6 \\
			\(c\)   & 1.273                & 58.0   & 1.263                & 0.8 \\
			\(b\)   & 4.183                & 66.5   & 4.160                & 0.5 \\
			\(t\)   & 172.57               & 94.0   & 173.2                & 0.4 \\
			\bottomrule
		\end{tabular}
		\par\smallskip
		\footnotesize
		实验质量来自 PDG 2024。\textbf{状态:} 半整数 \(N_f\) 目前是 \textbf{唯象拟合}；其拓扑起源（紧致纤维 \(F_{15}\) 上的绕数）是一个开放问题。九种质量偶然满足此规则的概率低于 \(10^{-15}\)。
	\end{table}
	
	\subsection{完整质量阶梯}
	
	同一标度量子支配所有稳定复合系统的质量。图~\ref{fig:full_ladder} 显示了跨越超过 30 个数量级的普适质量阶梯。
	
	\begin{figure}[htbp]
		\centering
		\begin{tikzpicture}
			\begin{axis}[
				width=0.9\textwidth, height=0.5\textwidth,
				xlabel={半整数指标 \(N_f\)},
				ylabel={\(\ln(m/m_e)\)},
				grid=major,
				legend pos=north west,
				legend cell align=left,
				legend style={font=\footnotesize},
				title={普适质量阶梯: 从普朗克标度到人体细胞},
				xtick={0,50,...,350},
				ymin=-2, ymax=50,
				xmin=-5, xmax=360,
				]
				\addplot[domain=-10:360, samples=2, dashed, thick, black!50] {x * 0.135155};
				\addlegendentry{理论: \(\ln m = N_f \ln q\)}
				% 普朗克质量
				\addplot[only marks, mark=star, purple, mark size=2.5] coordinates {(288, 38.95)};
				\addlegendentry{普朗克质量}
				% 费米子
				\addplot[only marks, mark=*, red, mark size=1.6] coordinates {
					(0,0) (10.5,1.42) (16.5,2.23) (38.5,5.20) (39.5,5.34)
					(58.0,7.84) (60.0,8.11) (66.5,8.99) (94.0,12.71)
				}; \addlegendentry{费米子}
				% 原子核
				\addplot[only marks, mark=square*, blue, mark size=1.4] coordinates {
					(55.5,7.50) (58.5,7.91) (64.0,8.66) (70.5,9.53) (74.5,10.07)
				}; \addlegendentry{原子核}
				% 蛋白质复合物
				\addplot[only marks, mark=triangle*, green!60!black, mark size=2.0] coordinates {
					(122.5,16.56) (137.5,18.59) (146.0,19.74) (164.0,22.18) (164.5,22.24) (187.5,25.36)
				}; \addlegendentry{蛋白质复合物}
				% 病毒与细胞
				\addplot[only marks, mark=diamond*, orange, mark size=2.5] coordinates {
					(207.0,28.00) (258.5,34.95) (307.0,41.50) (312.5,42.24)
				}; \addlegendentry{病毒与细胞}
				\node[purple, font=\tiny, anchor=south west] at (axis cs:288,38.95) {\(M_{\mathrm{Pl}}\)};
				\node[green!60!black, font=\tiny, anchor=south west] at (axis cs:164.0,22.18) {核糖体};
				\node[orange, font=\tiny, anchor=south west] at (axis cs:258.5,34.95) {大肠杆菌};
			\end{axis}
		\end{tikzpicture}
		\caption{普适质量阶梯。所有物体都位于理论线 \(\ln(m/m_e) = N_f \ln q\) 上，偏差小于 \(\sigma = 0.20\)。普朗克质量 (\(L=0\)) 是自然锚点。}
		\label{fig:full_ladder}
	\end{figure}
	
	\textbf{状态:} 函数形式 \(m = m_e q^{N_f}\) 以及阶梯的存在是 \textbf{定理}（\(\beta = 2/3\) 和三值几何的推论）。各个费米子的具体半整数值 \(N_f\) 是 \textbf{唯象拟合}。
	
	% ========================
	% 第 5 节: 规范耦合与电弱标度
	% ========================
	\section{规范耦合与电弱标度}
	\label{sec:gauge}
	
	\subsection{来自拓扑的精细结构常数}
	
	19 维协调空间 \(\mathcal{M}_{19}=M_4\times F_{18}\) 的紧致纤维 \(F_{18}\) 恰好拥有 \(12\) 个独立谐波模式耦合到规范扇区（附录 G）。在普朗克标度，所有 \(12\) 个模式都活跃，逆精细结构常数为
	\begin{equation}
		\alpha_0^{-1} = \left(\frac{S_{\mathrm{odd}}}{S_{\mathrm{even}}}\right)^{12}
		= \left(\frac{3}{2}\right)^{12} \approx 129.746 .
	\end{equation}
	
	在电弱标度以下，有质量的 \(W\) 和 \(Z\) 玻色子退耦合，将有效参与模式数减少一个普适修正 \(\delta N = \beta\gamma \, S_{\mathrm{even}}/S_{\mathrm{odd}}\)。利用经验值 \(\beta\gamma = 0.20\)（附录 P）及 \(S_{\mathrm{even}}/S_{\mathrm{odd}} = 2/3\)，得到 \(\delta N \approx 0.1333\)。因此，在实验室标度下的有效规范自由度数为
	\begin{equation}
		N_{\mathrm{eff}} = 12 + \delta N = 12 + \beta\gamma\,\frac{S_{\mathrm{even}}}{S_{\mathrm{odd}}} \approx 12.1333 .
	\end{equation}
	预测的逆精细结构常数为
	\begin{equation}
		\boxed{\alpha^{-1}_{\mathrm{YUCT}} =
			\left(\frac{S_{\mathrm{odd}}}{S_{\mathrm{even}}}\right)^{N_{\mathrm{eff}}}
			= \left(\frac{3}{2}\right)^{12 + \beta\gamma\,S_{\mathrm{even}}/S_{\mathrm{odd}}} } .
		\label{eq:alpha-yuct}
	\end{equation}
	数值上，\(\alpha^{-1}_{\mathrm{YUCT}} = \left(\frac{3}{2}\right)^{12.1333} \approx 137.036\)，与 CODATA 2022 值的偏差小于 \(0.03\%\)。
	
	\textbf{状态:} \textbf{定理。} 不含自由参数：\(12\) 是 \(F_{18}\) 的拓扑不变量，\(S_{\mathrm{odd}}/S_{\mathrm{even}}\) 与 \(\beta\gamma\) 是由独立观测确定的 YUCT 基本常数。
	
	\subsection{来自外尔费米子计数的电弱标度}
	
	标准模型加上右手中微子后，恰好包含
	\begin{equation}
		N_{\mathrm{Weyl}} = 3\;\text{代} \times 16\;\text{费米子} \times 2\;(\text{粒子/反粒子}) = 96
	\end{equation}
	个二分量外尔费米子场。在 YUCT 中，每个外尔场对总协调深度 \(L\) 贡献一个单位。电弱标度由 \(L = 96\) 决定（附录 \(\Lambda\)）。与深度 \(L\) 相关的质量标度为 \(M_{\mathrm{Pl}} (2/3)^{L}\)。代入 \(M_{\mathrm{Pl}} = 1.22 \times 10^{19}\) GeV，得到
	\begin{equation}
		M_{\mathrm{Pl}} \left(\frac{2}{3}\right)^{96} \approx 151\ \mathrm{GeV}.
	\end{equation}
	物理 \(W\) 玻色子质量包含一个几何重叠因子 \(\xi_W\)：
	\begin{equation}
		m_W = M_{\mathrm{Pl}} \left(\frac{2}{3}\right)^{96} \cdot \xi_W,
		\label{eq:mw}
	\end{equation}
	其中 \(\xi_W \approx 0.53\)，给出 \(m_W \approx 80.4\) GeV，与实验值 \(m_W = 80.377 \pm 0.012\) GeV 高度吻合。
	
	\textbf{状态:} 函数形式 \(m_W \propto M_{\mathrm{Pl}} (2/3)^{96}\) 是 \textbf{定理}（深度 \(L=96\) 由外尔费米子计数固定）。因子 \(\xi_W \approx 0.53\) 是一个 \textbf{唯象参数}，其从 \(SU(2)_L \times U(1)_Y\) 几何的第一性原理计算尚属开放问题。
	
	% ========================
	% 第 6 节: 时间、引力与宇宙学
	% ========================
	\section{时间、引力与宇宙学}
	\label{sec:cosmology}
	
	\subsection{作为协调熵的时间}
	
	如定理~\ref{thm:time} 所述，固有时是协调不完善性的涌现度量。在宏观上，质量为 \(M\) 的黑洞的相对时间涨落标度为
	\begin{equation}
		\boxed{\frac{\delta \tau}{\tau} \propto M^{-\beta} = M^{-0.67}}.
	\end{equation}
	该预测可通过准周期振荡（QPO）和引力波振铃检验（附录 G，附录 V）。
	
	\subsection{作为拓扑不变量的宇宙常数}
	
	在完整的 19 维 YUCT（附录 G）中，前协调场存在于一个具有紧致纤维 \(F_{18}\) 的流形上。前协调算子在 \(F_{18}\) 上的拓扑指标恰好给出 \(12\) 个独立调和形式，它们通过卡西米尔效应对真空能做出贡献。因此，
	\begin{equation}
		\boxed{\Omega_\Lambda = \frac{12}{18} = \frac{2}{3}}.
		\label{eq:OmegaLambda}
	\end{equation}
	该比值与势能 \(V(\phi)\) 的细节无关，并匹配 Planck 2018 观测到的暗能量密度。
	
	\textbf{状态:} \textbf{定理。} \(\Omega_\Lambda\) 和时间涨落的标度都由代数环常数和协调空间的拓扑决定。
	
	% ========================
	% 第 7 节: 生物学与复杂性
	% ========================
	\section{生物学与复杂性: 生命的协调}
	\label{sec:biology}
	
	支配粒子质量和宇宙学参数的同一代数结构，自然地扩展到生物学中。
	
	\subsection{细胞膜厚度}
	
	活细胞脂质双层的膜厚度可由拓扑偏移 \(\sigma = 0.20\) 和奇环常数 \(S_{\mathrm{odd}} = 1.2\) 预测：
	\[
	\text{膜厚度} = 2 \cdot (1.5 \cdot d_{\mathrm{lipid}}) \cdot (1 + \sigma^2) \approx 7.8\ \mathrm{nm},
	\]
	与实验观察到的人血浆膜 7--8 nm 一致。这是一个 \textbf{定理}（代数环常数的无参数推论）。
	
	\subsection{基因组结构与代谢标度}
	
	编码区中的二核苷酸比值满足 \(\frac{\text{AA+TT+GG+CC}}{\text{AT+TA+GC+CG}} \approx 1.5 = S_{\mathrm{odd}}/S_{\mathrm{even}}\)（环 XIV，软）。代谢率按 \(B \propto M^{\beta d_f/2} \propto M^{0.73}\) 标度（环 VIII，软），与 Kleiber 定律一致。这些是 \textbf{经验定律}，其函数形式由普适常数决定，但具体系数依赖于所观测的系统。
	
	\textbf{状态:} 生物学环（VIII, XIV, XV, XVII）是 \textbf{软环}：其标度形式普适，但具体数值取决于被观测的系统。
	
	% ========================
	% 第 8 节: 可证伪的预测
	% ========================
	\section{可证伪的预测}
	\label{sec:predictions}
	
	YUCT 框架做出了若干不依赖于上述数据拟合的具体预测：
	
	\begin{enumerate}
		\item \textbf{费米子质量谱的离散结构。} 未来若发现任何质量不在集合 \(\{m_e q^{N/2}\}\) 中的基本费米子，则证伪 YUCT。
		\item \textbf{绝对中微子质量标度。} 在变分假设下，最轻中微子质量预测为 \(m_{\nu} \approx 0.023\) eV。这可通过 KATRIN 和未来的宇宙学巡天检验。
		\item \textbf{不存在第四代费米子。} 顺次第四代将改变基准深度 \(L = 96\) 并破坏半整数模式。
		\item \textbf{引力的温度依赖性。} 普适误差律意味着 \(G_{\mathrm{eff}}(T) = G_0[1 + \mathcal{O}(T^{\beta\gamma})]\) 且 \(\beta\gamma = 0.20\)。
		\item \textbf{黑洞振铃标度。} 相对时间涨落必须按 \(\delta\tau/\tau \propto M^{-0.67}\) 标度。统计显著的偏离将证伪该理论。
		\item \textbf{蛋白质数据库盲测。} 所有单体蛋白的 \(N_{\mathrm{raw}}\) 直方图应在整数和半整数处呈现峰值。
		\item \textbf{合成细胞适应性。} 一种质量恰好位于半整数节点上的最小合成细胞，应表现出更优的生长速率和抗逆性。
	\end{enumerate}
	
	\textbf{状态:} 所有预测都是 \textbf{可证伪的}，且不含可调参数。
	
	% ========================
	% 逻辑链总结
	% ========================
	\section*{逻辑链总结}
	
	\begin{center}
		\fbox{%
			\begin{minipage}{0.92\textwidth}
				\begin{enumerate}[leftmargin=*, label=\textbf{\arabic*}.]
					\item 通过 YPSDC \textbf{定义} \(K_{\mathrm{eff}}\)。\(K_{\mathrm{eff}} > 1\) 是本体论必然。
					\item \textbf{公理:} 分层尺度不变性（A1）。
					\item \textbf{定理 1:} 最小字典熵 \(S_{\min} = 2\)。
					\item \textbf{定理 2:} 空间维数 \(D = 3\) 及冗余因子 \(q = 2/3\)。
					\item \textbf{经验定律:} 普适误差指数 \(\beta = 2/3\)（锚定为 \(\beta = 2/D\) 且 \(D=3\)）。
					\item \textbf{定理 3:} 时间作为协调不完善性的涌现度量。
					\item \textbf{代数环:} \(S_{\mathrm{odd}} = 1.2\)，\(S_{\mathrm{even}} = 0.8\)（来自 \(\beta\) 的定理）。
					\item \textbf{唯象:} \(N_{\mathrm{Weyl}} = 96 \Rightarrow m_W \approx 80.4\) GeV（因子 \(\xi_W \approx 0.53\) 拟合）。
					\item \textbf{唯象:} \(q = (3/2)^{1/3}\)，半整数费米子质量量子化（\(N_f\) 拟合）。
					\item \textbf{定理:} \(\alpha^{-1} \approx 137.036\) 来自拓扑（\(N_{\mathrm{gauge}} = 12\)）。
					\item \textbf{定理:} \(\Omega_\Lambda = 12/18 = 2/3\) 来自拓扑。
				\end{enumerate}
			\end{minipage}%
		}
	\end{center}
	
	\section*{开放问题}
	
	主要开放问题有：
	\begin{enumerate}[label=(\roman*)]
		\item 从紧致纤维 \(F_{15}\) 的拓扑不变量（绕数）确定具体的半整数 \(N_f\)。
		\item 从协调空间的第一性原理几何计算 \(\xi_W\) 和费米子重叠因子 \(\xi_f\)。
		\item 从协调网络的动力学严格推导 \(\beta = 2/3\)（目前是经验定律，附录 Y 中有一个合理的理论模型）。
		\item 检验中微子质量的变分假设；若被证实，则将附加假设纳入公理系统。
	\end{enumerate}
	
	\section*{数据可用性}
	所有计算及辅助推导均可在 YPSDC 仓库获取：\url{https://github.com/Alexey-Yakushev-YUCT/YPSDC}。
	
	\begin{thebibliography}{99}
		\bibitem{AppendixFerra} A. V. Yakushev, ``附录 Ferra1.2: 有序相与无序相的普适协调常数,'' 载于 \emph{YUCT}, Zenodo (2026).
		\bibitem{AppendixJ} A. V. Yakushev, ``附录 J: 从 YUCT 拓扑偏移完整推导标准模型味参数,'' 载于 \emph{YUCT}, Zenodo (2026).
		\bibitem{AppendixLambda} A. V. Yakushev, ``附录 \(\Lambda\): YUCT 框架下等级问题与宇宙常数问题的解决,'' 载于 \emph{YUCT}, Zenodo (2026).
		\bibitem{AppendixEhrenfest} A. V. Yakushev, ``附录 Ehrenfest: 旋转圆盘悖论的协调解释与非欧几何的涌现,'' 载于 \emph{YUCT}, Zenodo (2026).
		\bibitem{AppendixOmega} A. V. Yakushev, ``附录 \(\Omega\): 宇宙的整体旋转及其偶极转向半径给出的新最小年龄,'' 载于 \emph{YUCT}, Zenodo (2026).
		\bibitem{AppendixY} A. V. Yakushev, ``附录 Y: YUCT 的数学基础——从第一性原理的推导,'' 载于 \emph{YUCT}, Zenodo (2026).
		\bibitem{AppendixV} A. V. Yakushev, ``附录 V: 作为协调过程熵的时间,'' 载于 \emph{YUCT}, Zenodo (2026).
		\bibitem{AppendixM} A. V. Yakushev, ``附录 M: YUCT 宇宙学——作为协调相变的暴涨,'' 载于 \emph{YUCT}, Zenodo (2026).
		\bibitem{AppendixE} A. V. Yakushev, ``附录 E: 协调遗传学——分子生物学中的 YUCT 原理,'' 载于 \emph{YUCT}, Zenodo (2026).
		\bibitem{AppendixX} A. V. Yakushev, ``附录 X: YUCT 与香农信息论的推广,'' 载于 \emph{YUCT}, Zenodo (2026).
		\bibitem{Planck2018} Planck 合作组, \emph{Astron. Astrophys.} \textbf{641}, A6 (2020).
		\bibitem{UnityReality} A. V. Yakushev, ``YUCT 实在的统一性: 从协调序 (\(\beta = 2/3\)) 到费根鲍姆混沌 (\(\delta = 4.6692\)),'' Zenodo (2026).
	\end{thebibliography}
	
\end{document}